\chapter{Background and problem context}
\label{cap:background}
This chapter describes how the task of speech and music classification was first introduced to the public and how it evolved into a foundational building block of modern audio processing pipelines, summarizes the physical and acoustic characteristics that separate the two signals, and concludes with the definition and constraints of the on-line classifier.

\section{Origins and applications}
Automatic classification of speech and music originated in radio broadcasting, where it enabled receivers to automatically skip past speech and commercials in favor of music~\cite{saunders1996realtime}.

The growth of speech and music codecs introduced a parallel need for content classifiers at the codec front end: speech and music are best compressed by different codec families, and a misclassification at the input degrades compression quality~\cite{scheirer1997discriminator}.

Today, speech and music classification is a fundamental component of modern audio processing systems.
It is used in voice-activated assistants, music streaming platforms, and telecommunication systems, where accurate content identification enables appropriate codec selection and efficient resource usage~\cite{neuendorf2013usac}.
Beyond consumer applications, it also supports large-scale infrastructures such as broadcast monitoring and audio archiving.

\section{Characteristics of speech and music signals}
\label{sec:characteristics}
Speech differs from music along several physical and acoustic dimensions that every classifier in this work exploits.
The differences arise from how each signal is produced and shape the time-frequency patterns visible in their spectrograms.
The remainder of this section summarizes the properties that motivate the features and architectures introduced in \Cref{cap:classification_approaches}.

Speech is produced by a vocal tract whose shape changes slowly between voiced and unvoiced sound units, and the glottal source is quasi-periodic with a pitch span of roughly three octaves~\cite{khonglah2016gmmsvm}.
Music is produced by instruments whose tones span up to six octaves, decay slowly, and fall on a fixed set of frequencies set by the instrument's construction~\cite{bhattacharjee2018sps}.
The alternation of vowel and consonant in speech also produces large frame-to-frame energy variation that does not appear in sustained musical tones.

These production differences imprint distinctive patterns on the log-mel spectrogram.
Speech shows smooth arc-like striations whose harmonics shift continuously between sound units, while music shows long horizontal segments whose harmonics remain stationary before sharp transitions~\cite{bhattacharjee2018sps}.
Speech also concentrates modulation energy near the syllabic rate of about $4$\,Hz, a pattern that is largely absent in instrumental music~\cite{khonglah2016gmmsvm}.
Hand-engineered features and learned spectrogram representations both target these regularities.

\section{On-line classification}
\label{sec:online_classification}
This thesis will use the definitions of the term ``on-line processing'' according to Eyben's doctoral thesis~\cite{eyben2016realtime}.
On-line processing implicitly requires ``real-time processing'', so both terms are introduced.

\textit{``Real-time processing refers to processing of data in real-time.
This implies two things: (a) the processing of the data requires less time than the duration of the data, i.e., the time it will take to record or play-back the data at normal speed, and (b) the processing lag is as low as possible, ideally zero.''}~\cite[p.~139]{eyben2016realtime}

\textit{``On-line processing refers to processing of a continuous live stream of data while data are recorded or transmitted, i.e., computing and returning results continuously before the end of the data stream has been seen.''}~\cite[pp.~139--140]{eyben2016realtime}

With these definitions, on-line classification can be described as a process in which incoming data are processed continuously and classified into predefined classes, while maintaining low algorithmic delay.
While the definition allows small lookahead into future samples, I have decided to keep the models in this thesis strictly causal.

In practice, on-line processing usually requires the use of a buffer to aggregate incoming samples.
Once the buffer is filled, the classification is performed on the aggregated segment.
This aggregation introduces an \textbf{algorithmic delay} between data acquisition and the availability of a classification result: the minimum span of audio that must be received before the algorithm can produce any output.
For a causal system that operates on a fixed-length analysis frame (the unit of audio the algorithm processes to produce one decision), this cold-start delay equals the frame length and is independent of compute speed.
It would persist on an arbitrarily fast machine.
After the first frame, each subsequent decision waits only for one hop of new audio, since the prior analysis context is already buffered.
The steady-state algorithmic delay is therefore the hop length, typically in the range of 10--30\,ms, while the cold-start delay equals the frame length and is incurred once at stream start.
Features that aggregate statistics over a longer temporal context may rely on an additional rolling buffer of past samples.
This buffer carries prior context that has already been received and therefore does not contribute to algorithmic delay.
It is a context window, not a waiting period.

Algorithmic delay is only part of the picture.
A second contribution to the overall \textbf{latency} of an on-line classification system is the \textbf{computational delay}: the wall-clock time (the elapsed time on a real-world clock during execution) the implementation takes to extract features and produce a class label once an analysis frame is available.
Computational delay is hardware-dependent and grows with the complexity of the feature extractor and classifier.
Other contributions (audio acquisition, output dispatch, scheduling) are negligible compared to algorithmic and computational delay and are folded into the computational-delay budget for the rest of this thesis.
Total latency is therefore \textbf{algorithmic delay + computational delay}, with the first term fixed by the model definition and the second measured per model in \Cref{subsec:exp_complexity_results}.
On-line operation requires both to be small: algorithmic delay through a bounded analysis-frame duration, computational delay through computationally efficient algorithms with predictable execution times.

\subsection{Constraints}
\label{subsec:constraints}
Based on the above considerations, the following constraints apply to the on-line classification approach used in this work:
\begin{itemize}
    \item \textbf{Causality}: The classification must be performed using only past and present samples, without access to future data.
    \item \textbf{Real-time operation}: All processing steps must complete within the time duration of the processed data segment.
    \item \textbf{Low algorithmic delay}: The hop between consecutive decisions imposed by the algorithm must be kept short, since it lower-bounds the steady-state response time of the system.
    \item \textbf{Low computational delay}: Feature extraction and classification must have predictable and sufficiently low wall-clock cost, so that decisions are produced faster than new hops of audio arrive.
\end{itemize}


